%! TeX root: ../article.tex
% Anexo de resultados complementarios: contenido trasladado del cuerpo al
% recortar a ~20pp. Calibración E2/E3/E4, tablas/figuras de detalle y el
% detalle estadístico de la correlación interna-externa.

\section{Resultados complementarios}
\label{anexo:resultados}
Este anexo recoge material de resultados trasladado del cuerpo: el desarrollo
completo de los barridos de calibración E2/E3/E4, tablas y figuras de detalle
cuyo contenido se resume en el cuerpo, y el detalle estadístico de la
correlación interna--externa.

\subsection{Calibración: reducción de dimensionalidad}
\label{subsubsec:res-reduccion}
Reducir la dimensionalidad de los vectores de características es un paso común
en los \textit{pipelines} de agrupamiento. Sobre los vectores de 128
dimensiones de \texttt{dinov2\_graffiti\_style\_head}, el barrido (E2) compara
seis configuraciones y revela a UMAP como la única familia que produce
simultáneamente un número de clústeres no trivial, ruido moderado y tamaños
balanceados. La tabla~\ref{tab:e2-quality-6416} resume el resultado al tamaño
de dataset completo.

{\footnotesize
\begin{longtable}{|p{2.0cm}|p{1.0cm}|p{1.0cm}|p{1.2cm}|p{1.0cm}|p{1.0cm}|p{1.0cm}|} \hline
  \textbf{Reducción} & \textbf{cls} & \textbf{ruido} & \textbf{sil} & \textbf{CH} & \textbf{DB} & \textbf{csCV} \\ \hline
  \endhead
  identidad   & 301 & 0.684 &  \textbf{0.307} & 37.5 & \textbf{1.16} & \textbf{0.55} \\ \hline
  PCA-10      & 2   & 0.010 &  0.000 &  4.2 & 1.39 & 1.00 \\ \hline
  PCA-50      & 238 & 0.698 &  0.238 & 36.6 & 1.29 & 2.29 \\ \hline
  UMAP-10     & 562 & 0.202 &  0.155 & 27.9 & 1.70 & 0.74 \\ \hline
  UMAP-50     & 558 & \textbf{0.195} &  0.154 & 27.8 & 1.71 & 0.70 \\ \hline
  Isomap-10   & 125 & 0.624 &  0.044 & \textbf{38.4} & 1.80 & 3.40 \\ \hline
  \caption[E2, calidad intrínseca a $N=6416$]{E2, calidad intrínseca a $N=6416$. UMAP es la única reducción que
  produce simultáneamente un número de clústeres no trivial, ruido moderado y
  tamaños balanceados. \glosacalidad}
  \label{tab:e2-quality-6416}
\end{longtable}
}

La fila \emph{identidad} responde directamente a la pregunta de cuánto aporta
la reduccion de dimensionalidad: agrupando los 128-d crudos, HDBSCAN produce
particiones muy distintas según la escala. Colapsa a $k=2$ a $N=250$, pero a
$N=1000$ y $N=6416$ explota a 52 y 301 clústeres respectivamente. Esos
clústeres, sin embargo, se obtienen a costa de un ruido altísimo (0.605 y
0.684): HDBSCAN agrupa únicamente un núcleo denso y descarta entre el 60 y el
68\% de los puntos. La identidad no produce, por tanto, un colapso trivial,
pero tampoco una partición utilizable bajo el criterio conjunto, sino una
segmentación de cobertura muy baja. UMAP-10 sobre los mismos datos escala a 18,
106 y 562 clústeres manteniendo el ruido moderado (0.06, 0.11, 0.20) y los
tamaños balanceados. PCA-10 y PCA-50 ilustran otra inestabilidad: PCA-10
colapsa a $k=2$ con silueta nula ($0.000$) a $N=6416$ y desciende a silueta
negativa ($-0.020$) a $N=1000$, mientras que PCA-50 explota a $k=238$ con un
69.8\% de ruido y un coeficiente de variación de 2.29. Ambas son
configuraciones degeneradas que las métricas individuales premian
engañosamente: tanto la identidad (silueta $0.307$, DB $1.16$) como PCA-50
(silueta $0.238$) encabezan la tabla en silueta pese a descartar dos tercios
del dataset.

Las métricas de coste refuerzan la posición de la identidad como opción
inferior. Esta reducción no es gratuita como podría parecer, sino que empuja el
coste hacia la etapa de agrupamiento. Con $N=6416$ imágenes, HDBSCAN tarda
6.83~s sobre la representación identidad frente a 0.34~s sobre UMAP-10, unas
20 veces más. Isomap-10 fracasa en la misma vía, pero de forma más grave:
23.9~s y 1.25~GB de RSS pico, lo más caro del experimento, para una partición
igualmente degenerada (ruido 0.624 y csCV 3.40 a $N=6416$, con colapso a $k=2$ a
$N=1000$). UMAP-10 y UMAP-50 son indistinguibles en calidad, pero UMAP-50 cuesta
casi el doble en tiempo de reducción (11.0~s vs 5.9~s a $N=6416$). Por esto, la
recomendación es específicamente UMAP con \texttt{n\_components=10}, no UMAP en
general.

\begin{figure}[h!]
  \centering
  \includegraphics[width=\linewidth]{figures/e2_reduction_quality_cost.pdf}
  \caption[E2, calidad y coste por reducción]{E2: calidad intrínseca (panel izquierdo) y coste por etapa (panel
  derecho) por familia de reducción a $N=6416$. La fila identidad produce una
  partición de cobertura muy baja (ruido $\sim$0.68) y traslada el coste a la
  etapa de \textit{clustering}.}
  \label{fig:e2-reduction}
\end{figure}

\subsection{Calibración: algoritmo de agrupamiento}
\label{subsubsec:res-agrupamiento}
HDBSCAN es el único agrupador del panel que combina una silueta competitiva,
escalado sensato del número de clústeres con $N$, ruido moderado, tamaños
balanceados y coste despreciable bajo el criterio conjunto. El barrido (E3)
cubre 14 configuraciones sobre el extractor base y UMAP-10: HDBSCAN, K-medias
(ajustado con el método del codo), GMM diagonal (BIC con tres techos), DBSCAN
(auto-$\varepsilon$ y tres $\varepsilon$ fijos), OPTICS coseno, aglomerativo
(\textit{average linkage}, $k\in\{10, 20\}$) y espectral ($k\in\{10, 20\}$); la
tabla~\ref{tab:e3-summary-6416} extrae la fila representativa de cada familia.

{\footnotesize
\begin{longtable}{|p{2.4cm}|p{0.9cm}|p{1.0cm}|p{1.1cm}|p{0.8cm}|p{0.8cm}|p{0.9cm}|p{1.0cm}|p{1.0cm}|} \hline
  \textbf{Algoritmo} & \textbf{cls} & \textbf{ruido} & \textbf{sil} & \textbf{CH} & \textbf{DB} & \textbf{csCV} & \textbf{clu (s)} & \textbf{mem (MB)} \\ \hline
  \endhead
  HDBSCAN (mcs=5)        & 563 & 0.199 &  0.151 &  33 & 1.70 & 0.75 & 0.348 & 0.3 \\ \hline
  OPTICS (cos)           & 648 & 0.247 &  \textbf{0.177} &  30 & \textbf{1.64} & 0.45 & 8.578 & 763.3 \\ \hline
  GMM·300 (BIC)          & 297 & \textbf{0.000} &  0.054 &  50 & 2.52 & 0.58 & 1.839 & 49.7 \\ \hline
  K-medias (codo, $k=8$)   &   8 & \textbf{0.000} &  0.057 & \textbf{805} & 3.23 & \textbf{0.37} & \textbf{0.009} & \textbf{0.2} \\ \hline
  DBSCAN auto-$\varepsilon$ & 239 & 0.059 & -0.015 & 46 & 1.97 & 2.47 & 0.080 & \textbf{0.2} \\ \hline
  Aglom. \textit{avg} $k=10$ &  10 & \textbf{0.000} & -0.015 & 506 & 3.11 & 1.36 & 0.408 & 314.2 \\ \hline
  Espectral $k=10$       &  10 & \textbf{0.000} & -0.117 & 264 & 3.33 & 1.56 & 2.073 & 1.1 \\ \hline
  \caption[E3, fila representativa por familia a $N=6416$]{E3, fila representativa de cada familia a $N=6416$. La columna
  \textbf{mem} reporta el pico de RSS delta de la etapa de agrupamiento. Sólo
  HDBSCAN combina un número no trivial de clústeres, ruido moderado, tamaños
  balanceados, coste despreciable y huella de memoria mínima. \glosacalidad{}
  \textit{clu}=tiempo de agrupamiento, \textit{mem}=pico de RSS delta.}
  \label{tab:e3-summary-6416}
\end{longtable}
}

Como se anticipó en \ref{subsec:res-calibracion}, es necesario leer las
métricas intrínsecas  en conjunto con el número de clústeres, el ruido y el
csCV; nunca de manera aislada. Esto es especialmente importante aquí puesto que
los algoritmos basados en densidad (HDBSCAN, DBSCAN, OPTICS) son juzgados solo
sobre el subconjunto que agrupan, mientras que los particionales (K-medias, GMM,
aglomerativo, espectral) lo son sobre el conjunto de datos completo. La silueta
penaliza menos a un método que se abstiene de los puntos difíciles y CH premia
la compactación, sin importar la estructura.

Bajo este criterio, HDBSCAN parece ser el único algoritmo que produce
resultados defendibles en el rango de escalas evaluado: obtiene un ruido
moderado y un balance de tamaños razonable, una silueta competitiva y un coste
despreciable. OPTICS es su único par en calidad (superándolo ligeramente en
todas las escalas en silueta y sustancialmente en variación de tamaños), pero
se ve completamente dominado en coste, con un tiempo de ejecución $\sim$25
veces mayor y un uso de memoria $\sim$2500 veces mayor que HDBSCAN a $N=6416$.

Por otro lado, entre los métodos particionales, GMM con selección de $k$ por
BIC es la mejor alternativa si es necesaria una agrupación sin ruido. El resto
de métodos quedan dominados por baja calidad (DBSCAN es inestable, espectral
presenta la peor silueta del panel) o por el escalado de memoria.

\begin{figure}[h!]
  \centering
  \includegraphics[width=\linewidth]{figures/e3_pareto_quality_cost.pdf}
  \caption[E3, calidad frente a coste por familia]{E3: calidad frente a coste por familia a $N=6416$ (eje horizontal
  logarítmico). Cada familia se representa por su mejor variante en silueta
  (marcador grande); los puntos pequeños son el resto de variantes. La línea
  discontinua marca el frente de Pareto: HDBSCAN (rojo) domina a todas excepto
  a OPTICS, que paga $\sim$25 veces más coste por $\sim$0.02 más de silueta, y
K-medias, el algoritmo más barato pero con una silueta 0.094 puntos inferior a
la de HDBSCAN.}
  \label{fig:e3-pareto}
\end{figure}

\begin{figure}[h!]
  \centering
  \includegraphics[width=\linewidth]{figures/e3_scaling_with_n.pdf}
  \caption[E3, escalado por familia con $N$]{E3: escalado con el tamaño del banco $N$ por familia. Para cada $N$
  se selecciona la variante de mejor silueta dentro de la familia. El coste
  (izquierda) y la huella de memoria (centro) crecen varios órdenes de
  magnitud para OPTICS y aglomerativo mientras HDBSCAN se mantiene plano. El
  número de clústeres (derecha) crece de forma sensata para HDBSCAN, OPTICS y
  GMM, frente a los métodos con $k$ fijo (K-medias, aglomerativo, espectral) o
  inestables (DBSCAN).}
  \label{fig:e3-scaling}
\end{figure}

\subsection{Calibración: granularidad de HDBSCAN}
\label{subsubsec:res-mincluster}
Una vez escogido HDBSCAN como algoritmo de agrupamiento, el siguiente paso es
calibrar su hiperparámetro de granularidad \texttt{min\_cluster\_size} (mcs).
Este parámetro controla el número mínimo de puntos que debe tener un clúster
para ser considerado como tal, y por tanto afecta directamente al número de
clústeres, la fracción de ruido y el balance de tamaños. El barrido (E4) cubre
seis valores de \texttt{mcs} entre 5 y 200, como se puede ver en la
tabla~\ref{tab:e4-mcs-6416} a $N=6416$.

{\footnotesize
\begin{longtable}{|p{1.4cm}|p{1.0cm}|p{1.0cm}|p{1.0cm}|p{1.2cm}|p{1.0cm}|p{1.0cm}|} \hline
  \textbf{mcs} & \textbf{cls} & \textbf{ruido} & \textbf{sil} & \textbf{CH} & \textbf{DB} & \textbf{csCV} \\ \hline
  \endhead
    5 & 563 & 0.199 & 0.151 &   33 & \textbf{1.70} & \textbf{0.75} \\ \hline
   10 & 154 & 0.277 & 0.037 &   56 & 2.20 & 1.78 \\ \hline
   25 &  26 & 0.228 & 0.047 &  209 & 2.66 & 1.81 \\ \hline
   50 &   4 & \textbf{0.026} & \textbf{0.193} & \textbf{1404} & 2.01 & 0.86 \\ \hline
  100 &   2 & 0.008 & 0.269 & 2765 & 1.27 & 0.40 \\ \hline
  200 &   2 & 0.011 & 0.270 & 2776 & 1.27 & 0.40 \\ \hline
  \caption[E4, barrido de \texttt{min\_cluster\_size} a $N=6416$]{E4, barrido de \texttt{min\_cluster\_size} a $N=6416$. La silueta y
  el CH alcanzan su máximo en los valores degenerados ($k{=}2$); el criterio
  conjunto sitúa el óptimo en \texttt{mcs}=5. \glosacalidad}
  \label{tab:e4-mcs-6416}
\end{longtable}
}

Los valores más altos de \texttt{mcs} (100 y 200), aunque presentan las
métricas intrínsecas más altas con tasas de ruido y variación de tamaños bajas,
son claramente degenerados, puesto que solo producen dos clústeres masivos que
absorben casi todos los puntos. El caso de 50 es, aunque muy similar, algo más
defendible. El modelo utilizado para la generación de los \textit{embeddings}
(\texttt{dinov2\_graffiti\_style\_head}) fue entrenado con cuatro categorías de
estilo, por lo que esta segmentación a cuatro clústeres podría ser interpretada
como un resultado razonable.

Los siguientes dos valores, 25 y 10 respectivamente, presentan un número de
clústeres mayor, más segmentado, pero a costa de una caída considerable en
todas las métricas intrínsecas y un aumento de ruido y variación de tamaños. El
caso de \texttt{mcs}=5, presenta un número de clústeres aún mayor, pero con
menor ruido y métricas superiores a los dos casos anteriores (aunque no mejores
que los valores degenerados). Se escoge \texttt{mcs}=5 como valor base para el
resto del proyecto por una preferencia por un número de clústeres más alto y la
sospecha de que el caso para \texttt{mcs}=50 es un caso de agrupamiento
demasiado grueso. Esta decisión se ve reforzada en el experimento E8a, donde
las métricas supervisadas favorecen a la configuración escogida frente a la
otra.

\subsection{Tablas y figuras de detalle}
\label{anexo:detalle}
Material de detalle cuyo resumen se presenta en el cuerpo (sección
\ref{sec:results}): tablas de coste y figuras redundantes con su representación
preferida en el cuerpo, conservadas aquí por completitud y reproducibilidad.

% --- movido del cuerpo (Batch 4): coste E1, silueta E1, coste E7, ejemplos de clúster E8a ---

{\footnotesize
\begin{longtable}{|p{5.4cm}|p{1.7cm}|p{1.0cm}|p{1.4cm}|p{1.4cm}|} \hline
  \textbf{Extractor} & \textbf{ingesta (s)} & \textbf{ips} & \textbf{RSS pico (MB)} & \textbf{VRAM (MB)} \\ \hline
  \endhead
  \texttt{mobilenet\_v3}                  & \textbf{808.8}  & \textbf{7.93} & 366.0 & \textbf{31.8}  \\ \hline
  \texttt{mob.\_graffiti\_author\_head}   & 879.1  & 7.30 & 406.4 & 34.5  \\ \hline
  \texttt{inception\_v3}                  & 973.7  & 6.59 & 337.2 & 116.1 \\ \hline
  \texttt{vgg16}                          & 980.8  & 6.54 & 315.0 & 551.7 \\ \hline
  \texttt{resnet50}                       & 1080.0 & 5.94 & 336.0 & 117.8 \\ \hline
  \texttt{dinov2\_vits14}                 & 1215.7 & 5.28 & 209.6 & 98.0  \\ \hline
  \texttt{dinov2\_graffiti\_style\_head}  & 1245.0 & 5.15 & 286.2 & 99.7  \\ \hline
  \texttt{clip\_vit\_b32}                 & 1247.0 & 5.15 & \textbf{193.1} & 591.9 \\ \hline
  \texttt{yolon}                          & 1301.0 & 4.93 & 853.8 & 65.4  \\ \hline
  \texttt{yolom}                          & 1407.4 & 4.56 & 824.9 & 192.9 \\ \hline
  \caption[E1, coste de ingesta a $N=6416$.]{E1, coste de ingesta a $N=6416$. La cabeza añade un sobrecoste
  despreciable sobre el \textit{backbone}; los \textit{backbones} YOLO son los
  más caros simultáneamente en tiempo y RSS. Se muestran 10 de los 12
  extractores (las cabezas \texttt{mob.\_graffiti\_style\_head} y
  \texttt{dinov2\_graffiti\_author\_head} se omiten por brevedad; su coste cae
  dentro del rango de la cabeza correspondiente ya presente.)}
  \label{tab:e1-cost-6416}
\end{longtable}
}

\begin{figure}[h!]
  \centering
  \includegraphics[width=\linewidth]{figures/e1_silhouette_by_embedding.pdf}
  \caption[E1, silueta por extractor]{E1: silueta por extractor a tres tamaños
      de banco. Las CNN entrenadas en ImageNet encabezan la columna, pero la
  cabeza estilística de DINOv2 invierte silueta/CH (ver texto).}
  \label{fig:e1-silhouette}
\end{figure}

{\footnotesize
\begin{longtable}{|p{1.0cm}|p{1.6cm}|p{1.2cm}|p{2.0cm}|p{2.0cm}|p{1.4cm}|p{2.0cm}|} \hline
  \textbf{$N$} & \textbf{ingesta (s)} & \textbf{red (s)} & \textbf{HDBSCAN (s)} & \textbf{K-medias (s)} & \textbf{aglo (s)} & \textbf{búsqueda (s)} \\ \hline
  \endhead
   250 &   58.7 &  0.253 & 0.0038 & 0.0025 & 0.0019 &   3.66 \\ \hline
   500 &  115.7 &  0.603 & 0.0069 & 0.0027 & 0.0046 &  16.32 \\ \hline
  1000 &  231.6 &  1.514 & 0.0144 & 0.0033 & 0.0117 &  63.28 \\ \hline
  2000 &  455.9 &  4.542 & 0.1106 & 0.0047 & 0.0427 & 306.81 \\ \hline
  3500 &  776.0 & 11.870 & 0.1863 & 0.0066 & 0.1405 & 995.52 \\ \hline
  6416 & 1294.9 &  6.237 & 0.3580 & 0.0098 & 0.5988 & 3266.60 \\ \hline
  \caption[E7, tiempo por etapa frente a $N$]{E7, tiempo por etapa frente a $N$ (extracto). A $N=6416$ la búsqueda
  por similitud SQLite domina absolutamente, con la ingesta como segundo coste.
La caída de UMAP a $N=6416$ se discute en el texto.}
  \label{tab:e7-cost}
\end{longtable}
}

\begin{figure}[h!]
  \centering
  \includegraphics[width=0.85\linewidth]{figures/cluster_simple.pdf}\\[2pt]
  \includegraphics[width=0.85\linewidth]{figures/cluster_elaborate.pdf}\\[2pt]
  \includegraphics[width=0.85\linewidth]{figures/cluster_characters.pdf}
  \caption[Ejemplos de recortes por clúster HDBSCAN en E8a]{Ejemplos de recortes
  de cada uno de los tres clústeres de E8a (de arriba a abajo: sencillo,
  elaborado y \textit{characters}).}
  \label{fig:e8a-clusters}
\end{figure}
\subsection{Correlación interna--externa: detalle estadístico}
\label{anexo:correlacion}
Detalle numérico de la correlación de Pearson silueta--ARI resumida en
\ref{subsubsec:res-correlacion}.

Esta última subsección responde a si las métricas intrínsecas guardan alguna
relación clara con las extrínsecas. Sobre los resultados de E8a y E8b se calcula
la correlación de Pearson entre la silueta y el ARI a nivel de celda (cada celda
es una combinación de extractor, reducción y agrupador, agregada sobre sus
repeticiones).

En E8a (estilos, 4 clases) la correlación es positiva y significativa: $r =
+0.78$ ($n = 48$, $p < 10^{-7}$) sobre los agrupadores de $k$ fijo (K-medias,
aglomerativo y espectral con $k=4$), y $r = +0.55$ ($n = 89$, $p < 10^{-7}$) al
incluir las celdas HDBSCAN, de $k$ libre. En E8b (autoría, 87 clases) el signo
se invierte ($r = -0.54$ sobre las celdas con más de un clúster), pero este
valor apenas tiene peso práctico, ya que todo el rango de ARI en E8b vive en una
banda cercana al azar de modo que la correlación, sea cual sea su $r$ o su $p$
nominales, solo describe el orden de particiones que en el mayor de los casos
es aleatorio. No indica que la silueta descarte una estructura de autoría
recuperable, sino que tal estructura no existe a este nivel. La
figura~\ref{fig:e8-correlacion} representa ambas relaciones.

\begin{figure}[h!]
  \centering
  \includegraphics[width=\linewidth]{figures/e8_correlation.pdf}
  \caption[Silueta frente a ARI en E8a y E8b]{Silueta frente a ARI por ejecución en E8a (estilo) y E8b (autoría).
    Las rectas son ajustes por mínimos cuadrados sobre los agrupadores de $k$
    fijo (azul, continua) y sobre todas las celdas (gris, discontinua); en E8b se
    omiten porque el ARI no se despega del azar.}
  \label{fig:e8-correlacion}
\end{figure}
